\documentclass{abntexto-ufc}

\ufcsetup{
  font = @UFC_FONT@,
  code = listings,
  algorithms = algpseudocodex
}

\makeatletter
\newcommand{\ufcTextProbe}{%
  \typeout{UFC-TEXT-FAMILY=\f@family}%
  \typeout{UFC-TEXT-FONTSIZE=\f@size}%
}
\newcommand{\ufcCodeProbe}{%
  \typeout{UFC-CODE-FAMILY=\f@family}%
  \typeout{UFC-CODE-FONTSIZE=\f@size}%
}
\newcommand{\ufcAlgorithmProbe}{%
  \typeout{UFC-ALGORITHM-FAMILY=\f@family}%
  \typeout{UFC-ALGORITHM-FONTSIZE=\f@size}%
}
\makeatother

\begin{document}

\ufcTextProbe

\makeatletter
{\lst@basicstyle\ufcCodeProbe}
\makeatother

\lstset{numbers=left,numbersep=8pt}
\legend{code}{Código de validação tipográfica}
\ufcSource{Elaboração própria.}
\begin{ufclisting}[here]
UFC-CODE-GEOMETRY-MARKER
\end{ufclisting}

\legend{algorithm}{Algoritmo de validação tipográfica}
\ufcSource{Elaboração própria.}
\begin{ufcalgorithm}[here][1]
  \State \ufcAlgorithmProbe $x \gets 1$ \Comment{UFC-ALGORITHM-GEOMETRY-MARKER}
\end{ufcalgorithm}

\end{document}
